\section{Conclusion}
In this work, we proposed an autonomous agent framework, HaluAgent, which is capable of bilingual hallucination detection in Chinese and English. In our approach, we first curated a multi-functional toolbox to extend the LLM’s capability to detect hallucinations. Next, we designed a fine-grained three-stage detection framework along with memory mechanism, including sentence segmentation, tool selection and verification, and reflection.
Then, we leveraged existing datasets to synthesize detection trajectories by employing GPT-4 to execute the detection process following the HaluAgent framework.
Finally, we fine-tuned Baichuan2-Chat 7B and 13B on the synthesized trajectories. The HaluAgent models achieved notable improvements on both in-domain and out-of-domain datasets, with performance comparable to or even higher than GPT-4 without tool enhancements. 
Due to its high flexibility and adaptability, HaluAgent can effectively serve as a hallucination detector when human users interact with LLMs in real-world scenarios. In future work, we will extend our method to deal with more types of tools and hallucinations.